\chapter{Cumulative Lab: Put the Pieces Together}
\label{ch:cumulative-lab}

\begin{plainlanguage}
Now it is your turn to work through the project in one connected pass. Start
with a question, make sure you understand the tables, load the data, check what
happened, and explain what you found. The goal is not to move quickly. The goal
is to make each choice clear enough that you could walk someone else through
it.
\end{plainlanguage}

\section{Learning objectives}

By the end of the lab, you should be able to:
\begin{itemize}
  \item state an analytical question and identify its unit of analysis;
  \item explain the grain and key of every source table you use;
  \item establish a PostgreSQL environment without hiding configuration choices;
  \item audit, import, type, and validate the OULAD source data;
  \item write introductory SQL that answers bounded descriptive questions; and
  \item keep an evidence log that separates observed results from interpretation.
\end{itemize}

\section{Scenario and deliverables}

Imagine that a learning-support team wants an initial orientation to the data.
The team asks:

\begin{quote}
How many student--module attempts appear in each final-result category, how
often do students appear in more than one module presentation, and which source
quality conditions should constrain later analysis?
\end{quote}

The wording matters. This is descriptive, not causal. It asks what the recorded
data contain, not why a student passed, failed, or withdrew.

Your lab submission has five small deliverables:
\begin{enumerate}
  \item a one-page source and grain map;
  \item an environment record containing the database path you used;
  \item an import reconciliation table;
  \item a quality-gate note; and
  \item a six-query SQL portfolio with an evidence log.
\end{enumerate}

\begin{projectconnection}
\textbf{The lab's primary unit.}
The primary unit is a \emph{student--module presentation attempt}, represented
by one row in \projectfile{studentInfo.csv}. The same student may contribute
more than one attempt. A count of rows is therefore not automatically a count
of distinct people.
\end{projectconnection}

\section{Part A: Write the source and grain map}

Before opening SQL, complete \cref{tab:lab-source-map}. The first two rows are
provided as examples.

\begin{table}[htbp]
\centering
\small
\caption{Cumulative lab source and grain map}
\label{tab:lab-source-map}
\begin{tabularx}{\textwidth}{@{}lXXX@{}}
\toprule
Source & One row represents & Candidate key & Intended role \\
\midrule
\projectfile{studentInfo.csv} &
One student in one module presentation &
\texttt{student\_id + module + presentation} &
Attempt-level demographics and final result \\
\projectfile{studentRegistration.csv} &
One registration record for an attempt &
\texttt{student\_id + module + presentation} &
Registration and withdrawal timing \\
\projectfile{courses.csv} & \emph{Complete} & \emph{Complete} & \emph{Complete} \\
\projectfile{assessments.csv} & \emph{Complete} & \emph{Complete} & \emph{Complete} \\
\projectfile{studentAssessment.csv} & \emph{Complete} & \emph{Complete} & \emph{Complete} \\
\projectfile{vle.csv} & \emph{Complete} & \emph{Complete} & \emph{Complete} \\
\projectfile{studentVle.csv} & \emph{Complete} & \emph{Complete} & \emph{Complete} \\
\bottomrule
\end{tabularx}
\end{table}

\begin{validationbox}
For every table, say the phrase ``one row represents'' aloud. If you cannot
finish the sentence precisely, return to \cref{ch:oulad-data}. A list of column
names is not a grain statement.
\end{validationbox}

\begin{exercisebox}{Lab checkpoint 1: defend the unit of analysis}
\label{ex:lab-grain}
Write three sentences explaining why \texttt{COUNT(*)} from
\texttt{core.student\_info} counts attempts, why
\texttt{COUNT(DISTINCT student\_id)} counts people, and why the two values need
not match. Then predict which count will be larger in this dataset.
\end{exercisebox}

\section{Part B: Record the environment}

Choose one of the supported paths from \cref{ch:postgresql-setup}: Docker
Compose or a local PostgreSQL installation. Record:
\begin{itemize}
  \item the path chosen;
  \item the PostgreSQL version;
  \item the database name;
  \item the exact form of \texttt{DATABASE\_URL}, with any password redacted;
  \item the command used to prove the connection; and
  \item the date and time of the run.
\end{itemize}

\begin{lstlisting}[language=bash,caption={A concise environment record}]
python --version
psql --version
psql "$DATABASE_URL" -X -v ON_ERROR_STOP=1 \
  -c "select current_database(), current_user, version();"
\end{lstlisting}

\begin{commonmistakes}
\textbf{Do not publish secrets.}
Record the \emph{shape} of a connection string, but redact passwords and private
host details before publishing screenshots, logs, or coursework.
\end{commonmistakes}

\begin{exercisebox}{Lab checkpoint 2: diagnose before proceeding}
\label{ex:lab-environment}
Suppose \texttt{psql} reports that the role
\texttt{learning\_analytics} does not exist. Give one repair for the Docker path
and one repair for a trusted local PostgreSQL path. Explain why repeatedly
rerunning the import script would not solve this connection error.
\end{exercisebox}

\section{Part C: Audit and import}

Follow the repository's intended sequence:
\begin{enumerate}
  \item obtain the source archive through the audited download command;
  \item verify the expected files and source checksum;
  \item initialize the schemas and raw tables;
  \item load all seven CSV files;
  \item build the staging and core relations; and
  \item reconcile source, raw, and core counts.
\end{enumerate}

Do not skip the source audit merely because the filenames look right. A
reproducible source is identified by content and structure, not only by a
familiar name.

\Needspace{13\baselineskip}
\begin{lstlisting}[language=bash,caption={Repository-level orchestration for a complete first run}]
make setup
make data
make db-init
make load
make pipeline
make quality
\end{lstlisting}

The exact targets may be run separately during diagnosis. The shell scripts use
stop-on-error behavior, and the SQL entry points use
\texttt{ON\_ERROR\_STOP}; a failed statement should interrupt the run rather
than quietly producing a partly built database.

\subsection{Reconciliation table}

Complete one row per source:

\begin{table}[htbp]
\centering
\small
\caption{Import reconciliation record}
\label{tab:lab-reconciliation}
\begin{tabularx}{\textwidth}{@{}lrrrrX@{}}
\toprule
Source & CSV rows & Raw rows & Core rows & Difference & Explanation \\
\midrule
\projectfile{courses.csv} & 22 & & & & \\
\projectfile{assessments.csv} & 206 & & & & \\
\projectfile{vle.csv} & 6,364 & & & & \\
\projectfile{studentInfo.csv} & 32,593 & & & & \\
\projectfile{studentRegistration.csv} & 32,593 & & & & \\
\projectfile{studentAssessment.csv} & 173,912 & & & & \\
\projectfile{studentVle.csv} & 10,655,280 & & & & \\
\bottomrule
\end{tabularx}
\end{table}

For the current source, the seven CSV files contain 10,900,970 data rows in
total. Treat that number as a reconciliation target, not as proof that every
row is analytically valid.

\begin{exercisebox}{Lab checkpoint 3: investigate a discrepancy}
\label{ex:lab-reconcile}
Assume \texttt{raw.student\_assessment} has 173,911 rows after import, while the
audited CSV has 173,912. List a three-step investigation that preserves the
failed evidence. Do not propose manually inserting a guessed row.
\end{exercisebox}

\section{Part D: Interpret the quality gate}

Run the repository quality suite and classify the output:
\begin{description}
  \item[Errors] Conditions that must pass before analysis continues.
  \item[Warnings] Conditions that require interpretation or a documented choice.
  \item[Information] Profiles that describe the data without declaring failure.
\end{description}

The checked project results contain 24 quality entries: 18 error-level tests
that pass, three warning profiles, and three informational profiles. The
warnings identify 688,988 activity records dated before module start, 11
assessments with missing due dates, and 1,111 attempt records with missing
IMD-band values.

\begin{exercisebox}{Lab checkpoint 4: write a quality-gate note}
\label{ex:lab-quality}
Write a five-sentence quality-gate note. It must state whether the error gate
passed, describe all three warning profiles, distinguish a warning from a
failure, and name one later analytical decision affected by each warning.
\end{exercisebox}

\section{Part E: Build the six-query portfolio}

Run each query, save the SQL, and add one evidence-log entry. The queries are
small by design: the challenge is precise interpretation.

\subsection{Query 1: confirm the attempt population}

\begin{lstlisting}[language=SQL]
SELECT
    COUNT(*) AS attempt_count,
    COUNT(DISTINCT student_id) AS student_count
FROM core.student_info;
\end{lstlisting}

Expected values for the audited source are 32,593 attempts and 28,785 distinct
students.

\subsection{Query 2: describe final outcomes}

\begin{lstlisting}[language=SQL]
SELECT
    final_result,
    COUNT(*) AS attempts,
    ROUND(100.0 * COUNT(*) / SUM(COUNT(*)) OVER (), 1) AS pct
FROM core.student_info
GROUP BY final_result
ORDER BY attempts DESC;
\end{lstlisting}

The expected counts are Pass 12,361, Withdrawn 10,156, Fail 7,052, and
Distinction 3,024. State explicitly that these are attempt counts.

\subsection{Query 3: count repeat appearances}

\begin{lstlisting}[language=SQL]
SELECT
    COUNT(*) FILTER (WHERE attempt_count = 1) AS one_attempt,
    COUNT(*) FILTER (WHERE attempt_count > 1) AS multiple_attempts,
    MAX(attempt_count) AS maximum_attempts
FROM (
    SELECT student_id, COUNT(*) AS attempt_count
    FROM core.student_info
    GROUP BY student_id
) AS per_student;
\end{lstlisting}

\subsection{Query 4: inspect withdrawal timing}

\begin{lstlisting}[language=SQL]
SELECT
    COUNT(*) AS registration_records,
    COUNT(date_unregistration) AS recorded_withdrawal_dates,
    MIN(date_unregistration) AS earliest_relative_day,
    MAX(date_unregistration) AS latest_relative_day
FROM core.student_registration;
\end{lstlisting}

The day fields are relative to presentation start. Do not silently translate
them into calendar dates.

\subsection{Query 5: profile assessment due-date missingness}

\begin{lstlisting}[language=SQL]
SELECT
    assessment_type,
    COUNT(*) AS assessments,
    COUNT(*) FILTER (WHERE due_date IS NULL) AS missing_due_dates
FROM core.assessments
GROUP BY assessment_type
ORDER BY assessment_type;
\end{lstlisting}

\subsection{Query 6: bound engagement timing}

\begin{lstlisting}[language=SQL]
SELECT
    COUNT(*) AS activity_rows,
    COUNT(*) FILTER (WHERE activity_date < 0) AS pre_start_rows,
    MIN(activity_date) AS earliest_day,
    MAX(activity_date) AS latest_day
FROM core.student_vle;
\end{lstlisting}

\begin{commonmistakes}
\textbf{A row is not a behavior.}
An activity record reflects the structure of the VLE extract. It is not,
without further definition, a learning session, a moment of attention, or
evidence of motivation.
\end{commonmistakes}

\section{The evidence log}

Use one row per reported finding:

\begin{table}[htbp]
\centering
\small
\caption{Evidence-log template}
\label{tab:evidence-log}
\begin{tabularx}{\textwidth}{@{}p{0.08\textwidth}XXXX@{}}
\toprule
Query & Observed result & Unit and population & Interpretation & Limitation or check \\
\midrule
1 & & & & \\
2 & & & & \\
3 & & & & \\
4 & & & & \\
5 & & & & \\
6 & & & & \\
\bottomrule
\end{tabularx}
\end{table}

A strong entry for Query 2 might read: ``Among 32,593 recorded
student--module presentation attempts, 10,156 (31.2\%) have a final result of
Withdrawn. This describes recorded outcomes at attempt grain; it does not
estimate a causal effect or a person-level withdrawal rate.''

\begin{exercisebox}{Lab checkpoint 5: final synthesis}
\label{ex:lab-synthesis}
Write a 150--200 word briefing for the learning-support team. Include two
descriptive findings, one quality warning, one grain warning, and one next
question. Avoid the words \emph{caused}, \emph{proved}, and \emph{at risk}
unless the evidence genuinely supports them.
\end{exercisebox}

\section{Common mistakes}

\begin{itemize}
  \item Reporting an attempt count as a student count.
  \item Treating an expected row count as sufficient evidence of source integrity.
  \item Continuing after a failed error-level quality check.
  \item Deleting pre-start activity because its date looks surprising.
  \item Calling missing values ``bad data'' before studying their meaning.
  \item Describing a descriptive association as an explanation or intervention effect.
  \item Publishing credentials or machine-specific paths in an evidence log.
\end{itemize}

\section{Transfer task}

Choose a different relational dataset with at least three tables. Produce only
the source-and-grain map and the six fields of an environment record. Then
write two introductory aggregation queries. The point is to transfer the
workflow, not the OULAD vocabulary.

\section{Guided solutions}

Guided responses for the five checkpoints appear in
\cref{sol:lab-grain,sol:lab-environment,sol:lab-reconcile,sol:lab-quality,sol:lab-synthesis}.

\section{Chapter summary}

The cumulative lab turns a collection of commands into an analytical
discipline: define the question, name the grain, record the environment, audit
the source, reconcile the import, pass the quality gate, query deliberately,
and preserve the boundary between evidence and interpretation.
